\section*{Exact stationary diffusion criterion}
Assume $J_m(a,b,H)$ is Hurwitz on the stoichiometric subspace and define
\[
\Theta(H,D)=d_Z-8h_Z\sum_{j=2}^{m-1}\frac{d_j}{h_j}.
\]
Then
\[
\boxed{\text{a nonzero stationary zero crossing occurs on the ray }sD
\iff \Theta(H,D)>0.}
\]
Because $c^TJ=0$ and $J|_{c^\perp}$ is Hurwitz, $J$ has rank $n-1$, its
conservation zero is algebraically simple, and
\[
 \sum_{|I|=n-1}(-1)^{|I|}\det J_{I,I}>0.
\]
Thus the principal-minor diffusion-ray theorem applies.  When $\Theta>0$,
the threshold $s_*(a,b,H,D)$ is unique and simple.  Every
$0<s<s_*(a,b,H,D)$ has a positive real eigenvalue.  Scaling $D$ by
$\gamma>0$ gives
\[
 s_*(a,b,H,\gamma D)=\gamma^{-1}s_*(a,b,H,D).
\]
When $\Theta\le0$, no
zero-eigenvalue stationary crossing occurs on the positive ray.  The statement
does not classify possible wave instability unless a separate mode theorem is
available.
